\chapter{Troponin}

\section*{What Is This Test?}
Cardiac troponin tests measure cardiac-specific isoforms of the troponin complex (Troponin I or T) released into the circulation when cardiomyocytes are injured. Together with clinical presentation and ECG changes, troponin is the cornerstone for diagnosis of acute myocardial infarction (universal definition of MI, Fourth Universal Definition, 2018). Troponin has superseded CK-MB and myoglobin because of superior sensitivity and specificity for myocardial injury. High-sensitivity troponin assays can detect concentrations at the 99th percentile of a healthy population with $\leq$ 10\% coefficient of variation, enabling early rule-in/rule-out protocols within 1--2 hours.

\section*{Alternative Names}
Troponin I (cTnI); Troponin T (cTnT); High-Sensitivity Troponin (hs-cTnI, hs-cTnT); Cardiac Troponin; hs-Troponin

\section*{Principle}
Sandwich immunoassay with cardiac-specific monoclonal antibodies capturing the troponin I or T molecule using two different epitopes. Signal generation by chemiluminescence (most analyzers: Roche Elecsys, Abbott Architect, Siemens ADVIA Centaur, Beckman Access) read by photomultiplier. High-sensitivity assays (hs-cTn) measure reliably at $\leq$ 3 ng/L with $\leq$ 10\% CV; standard assays require $>$ 50--100 ng/L for $\leq$ 10\% CV.

\section*{History}
Troponin was discovered as part of the contractile machinery of striated muscle in the 1960s. Cardiac-specific isoforms were identified in the 1980s and commercial immunoassays introduced in the 1990s. The rise of troponin as the preferred biomarker accelerated through the 2000s; the First Universal Definition of MI (2007) endorsed troponin superiority. High-sensitivity assays (Roche hs-cTnT, Abbott hs-cTnI) appeared in 2010s and became routine in the U.S. by 2017.

\section*{Why This Test Is Ordered}
\begin{itemize}
  \item Diagnosis of acute myocardial infarction in patients presenting with chest pain or ischemic equivalent symptoms
  \item Risk stratification and decision-making in acute coronary syndromes (STEMI, NSTEMI)
  \item Detection of myocardial injury in non-ACS conditions: sepsis, pulmonary embolism, takotsubo, severe heart failure, renal failure
  \item Monitoring of periprocedural myocardial injury after PCI, CABG, ablation
  \item Screening for cardiotoxicity in chemotherapy (anthracyclines, trastuzumab)
\end{itemize}

\section*{Is This a Primary or Secondary Test}
Primary biomarker for suspected MI; serial measurements required for diagnosis when baseline troponin is elevated or unclear.

\section*{How This Test Is Performed}
Venous blood drawn in serum separator or heparin tubes; rapid turnaround (bedside POC or central laboratory) within 30--60 minutes. For suspected ACS, draw at presentation and repeat at 1, 2, and 3--6 hours (algorithm-dependent). Stable patients with chronic elevations (e.g., renal failure) require delta values rather than absolute cutoffs.

\section*{What Preparation Is Needed}
None in emergency setting. Document time of symptom onset (threshold of troponin rise begins at 2--4 hours).

\section*{Contraindications and Precautions}
Troponin is cardiac-specific but not disease-specific: elevations occur in any condition that injures cardiomyocytes, including sepsis, renal failure, PE, stroke, stress cardiomyopathy, myocarditis, heart failure. Persistent chronic elevations in renal failure require interpretation with delta values. Heterophilic antibodies may falsely elevate troponin I (especially cTnI assays); confirmatory testing on different platforms if suspicion.

\section*{Risks}
Venipuncture only.

\section*{What Happens After the Test}
Result reviewed with ECG and clinical context. Rising pattern (delta) confirms type 1 MI (plaque rupture) or type 2 MI (supply-demand mismatch). Persistent elevation without delta may reflect chronic myocardial injury (renal failure, stable heart failure).

\section*{Understanding Results}

\subsection*{Below Normal}
A normal troponin with chest pain onset $>$ 3 hours ago has high negative predictive value for AMI; rapid rule-out protocols (ESC 0/1, 0/1-h) reduce unnecessary admissions. Very low or undetectable high-sensitivity troponin rules out MI in emergency departments.

\subsection*{Above Normal}
\begin{itemize}
  \item \textbf{Type 1 MI (plaque rupture, STEMI/NSTEMI):} Rise and/or fall of troponin with $\geq$ 1 value $>$ 99th percentile URL, plus ischemic symptoms, ECG changes, or imaging evidence
  \item \textbf{Type 2 MI (supply-demand mismatch):} Elevated troponin without plaque rupture (tachyarrhythmia, hypotension, severe anemia)
  \item \textbf{Acute myocardial injury (non-MI)}: Sepsis, myocarditis, takotsubo, pulmonary embolism (right ventricular strain), acute heart failure
  \item \textbf{Chronic myocardial injury}: Persistently elevated troponin without delta; structural heart disease, renal failure, advanced heart failure
  \item \textbf{Chemotherapy cardiotoxicity}: Persistent elevation after anthracycline therapy
\end{itemize}

\section*{Normal Results}
\begin{itemize}
  \item \textbf{Standard troponin I/T:} $<$ 0.04--0.04 ng/mL (assay-dependent)
  \item \textbf{High-sensitivity troponin I (Abbott):} 99th percentile URL $<$ 16 ng/L (women), $<$ 34 ng/L (men)
  \item \textbf{High-sensitivity troponin T (Roche):} 99th percentile URL $<$ 14 ng/L
  \item \textbf{MI definition:} Troponin $>$ 99th percentile URL with significant delta rise/fall
  \item \textbf{ESC 0/1-h algorithm delta thresholds}: Based on absolute change at 1 hour (assay-specific)
\end{itemize}

\section*{Follow-up Tests}
\begin{itemize}
  \item \textbf{ECG (serial):} ST elevation, T-wave changes, bundle branch block
  \item \textbf{Echocardiography (urgent):} Regional wall motion abnormality in acute MI; right ventricular strain in PE; apical ballooning in takotsubo
  \item \textbf{ Serial troponin at 1, 2, and 3--6 hours}: Define delta pattern
  \item \textbf{Coronary angiography}: For STEMI or unstable NSTEMI patients
  \item \textbf{BNP/NT-proBNP}: Concurrent heart failure assessment
  \item \textbf{D-dimer, CT pulmonary angiography}: If PE suspected
  \item \textbf{CK-MB (historic; rarely needed)}: When troponin ambiguous or suspected reinfarction
\end{itemize}

\section*{References}
\begin{enumerate}
  \item Thygesen K, Alpert JS, Jaffe AS, et al. Fourth Universal Definition of Myocardial Infarction. J Am Coll Cardiol. 2018;72(18):2231-2264.
  \item Collet JP, Thiele H, Barbato E, et al. 2020 ESC guidelines for the management of acute coronary syndromes. Eur Heart J. 2021;42(14):1289-1367.
  \item Shah ASV, Anand A, Sandoval Y, et al. High-sensitivity cardiac troponin I at presentation. Lancet. 2015;386(10012):2481-2488.
  \item Morrow DA, et al. APACE-Troponin T risk score. Circulation. 2014;129(1):56-65.
\end{enumerate}

\clearpage